% ====================================================
% الملحق O – القياس العالمي للأحجام الحرجة للمجموعات
% ====================================================
\documentclass[12pt]{article}
\usepackage[utf8]{inputenc}
\usepackage[T1]{fontenc}
\usepackage[english]{babel}   % للغة الإنجليزية في بعض الأماكن

% ============================================================
% إعدادات اللغة العربية (LuaLaTeX / XeLaTeX)
% ============================================================
\usepackage{polyglossia}
\setmainlanguage{arabic}
\setotherlanguage{english}

% خطوط عربية – Amiri متوفر في النظام
\newfontfamily\arabicfont{Amiri}[Script=Arabic, Language=Arabic]
\newfontfamily\arabicfontsf{Amiri}[Script=Arabic, Language=Arabic]
\newfontfamily\arabicfonttt{Amiri}[Script=Arabic, Language=Arabic]

% خطوط لاتينية
\newfontfamily\englishfont{Times New Roman}[Ligatures=TeX]
\setmonofont{Latin Modern Mono}

% ============================================================
% حزم الرياضيات والتنسيق
% ============================================================
\usepackage{amsmath,amssymb,amsthm,mathtools}
\usepackage{geometry}
\usepackage{booktabs}
\usepackage{hyperref}
\usepackage{url}
\usepackage{tabularx}
\usepackage{array}
\usepackage{makecell}
\usepackage{ragged2e}
\usepackage{bm}
\usepackage{graphicx}
\usepackage{caption}
\usepackage{subcaption}
\usepackage{float}

% بيئات النظريات – عربية
\newtheorem{theorem}{نظرية}
\newtheorem{lemma}{ليما}
\newtheorem{definition}{تعريف}
\newtheorem{corollary}{نتيجة طبيعية}
\theoremstyle{remark}
\newtheorem{remark}{ملاحظة}

% أوامر مختصرة
\newcommand{\Keff}{K_{\mathrm{eff}}}
\newcommand{\eps}{\varepsilon}
\newcommand{\df}{d_{\!f}}
\newcommand{\kappac}{\kappa}
\newcommand{\betaf}{\beta}
\newcommand{\Tmin}{T_{\min}}
\newcommand{\Dprior}{D_{\mathrm{prior}}}

\hypersetup{colorlinks=true,linkcolor=blue,citecolor=blue,urlcolor=blue}

\begin{document}
	
	\begin{titlepage}
		\begin{center}
			\vspace*{0.5cm}
			\Huge\textbf{الملحق O. القياس العالمي للأحجام الحرجة للمجموعات: \\ أدلة تجريبية وارتباطها ببديهية العتبة التكنولوجية (YUCT)}
			
			\vspace{0.3cm}
			\small\textit{مجموعة منهجية من البيانات حول الأحجام الحرجة للمجموعات في الأنظمة البيولوجية والاجتماعية، مفسَّرة من خلال بديهية عتبة القاموس التكنولوجي (الملحق O) والأس العالمي \(\beta = 0.67\) (الملحق L).}
			
			\vspace{0.3cm}
			\small\textbf{أليكسي ف. ياكوشيف} \\
			نظرية YUCT: \url{https://yuct.org/}\\
			بروتوكول YPSDC: \url{https://ypsdc.com/}\\
			
			\vspace{2cm}
			\textbf{DOI:} \url{https://doi.org/10.5281/zenodo.18444598}\\
			
			\vspace{1cm}
			\large 
			نُشرت نسخة مختصرة من هذا العمل سابقًا وهي متاحة على\\
			\url{https://doi.org/10.5281/zenodo.18362308}\\
			\vspace{1cm}
			مارس 2026 \\ \large الإصدار 2.0.
			
			\vspace{0.5cm}
			\begin{minipage}{0.95\textwidth}
				\begin{abstract}
					\noindent
					يجمع هذا الملحق بيانات تجريبية حول الأحجام الحرجة التي تشهد عندها الأنظمة الجماعية المختلفة تحولاً طورياً – مثل الأسراب، الانشطار، أو إعادة التنظيم الهيكلي. تشمل الأنظمة المدروسة مستعمرات نحل العسل، الوحدات العسكرية، أسراب الأسماك، مجموعات الدلافين، تجمعات الماشية، ومجموعات الرئيسيات. باستخدام قانون الخطأ العالمي \(\varepsilon \propto N^{\beta}\) مع \(\beta = 0.67\) (الملحق L) وبديهية عتبة القاموس التكنولوجي (الملحق O)، نشتق علاقة كمية بين الحجم الحرج \(N_{\text{crit}}\) الذي يحدث عنده التحول والحجم الأمثل للمجموعة \(N_{\text{opt}}\). تتجمع النسب الملاحظة \(N_{\text{crit}}/N_{\text{opt}}\) للأنظمة المتنوعة حول قيم متسقة مع \(\beta = 0.67\)، مما يوفر دعمًا مستقلاً لعالمية هذا الأس وللفكرة القائلة بأن تجاوز عتبة تكنولوجية (أو تنظيمية) يقابل تحولاً طورياً في كفاءة التنسيق. نقترح أيضًا تنبؤات محددة قابلة للاختبار لتجارب ودراسات ميدانية مستقبلية يمكن أن تزيد من التحقق من إطار YUCT.
				\end{abstract}
			\end{minipage}
			
			\vspace{0.3cm}
			\noindent
			\small\textbf{الكلمات المفتاحية:} YUCT، الحجم الحرج للمجموعة، التحول الطوري، سرب نحل العسل، التسلسل الهرمي العسكري، أسراب الأسماك، القياس العالمي، \(\beta = 0.67\)، العتبة التكنولوجية.
		\end{center}
	\end{titlepage}
	
	\tableofcontents
	\newpage
	
	\section{مقدمة}
	\label{sec:intro}
	
	أحد التنبؤات المركزية لنظرية ياكوشيف الموحدة للتنسيق (YUCT) هو وجود أس عالمي \(\beta = 0.67\) يحكم قياس أخطاء التنسيق مع حجم النظام (الملحق L). ينبثق هذا الأس من البنية الكسيرية للقواميس والقيود المعلوماتية، وقد تم التحقق منه عبر 40 مرتبة من حيث الحجم – من تضاعف الحمض النووي إلى تموجات الخلفية الكونية الميكروية. إذا كان \(\beta\) عالميًا حقًا، فينبغي أن يتجلى أيضًا في الأحجام الحرجة التي تخضع عندها الأنظمة الجماعية لتحولات طورية، مثل الأسراب، الانشطار، أو إعادة التنظيم الهرمي. علاوة على ذلك، تنص بديهية عتبة القاموس التكنولوجي (الملحق O) على أن النظام يجب أن يصل إلى مستوى تكنولوجي أدنى \(\Tmin\) ليتمكن من استخدام أو توليد قاموس؛ ويمكن إعادة تفسير هذه العتبة كنقطة حرجة في كفاءة تنسيق النظام \(\Keff\). يجمع هذا الملحق بيانات تجريبية من مجالات متنوعة ويظهر أن النسب الملاحظة للأحجام الحرجة إلى الأحجام المثلى تتفق مع تنبؤات YUCT، مما يربط البديهية المجردة بكميات ملموسة قابلة للقياس.
	
	\section{الإطار النظري}
	\label{sec:theory}
	
	بالنسبة لأي نظام منسق حجمه \(N\)، يُفترض أن خطأ التنسيق النسبي \(\varepsilon\) يخضع للقياس
	\begin{equation}
		\varepsilon(N) = \varepsilon_0 N^{\beta}, \qquad \beta = 0.67,
		\label{eq:scaling}
	\end{equation}
	حيث \(\varepsilon_0\) ثابت يعتمد على النظام. يعمل النظام على النحو الأمثل عند حجم \(N_{\text{opt}}\) حيث يكون \(\varepsilon\) في حده الأدنى. كلما زاد \(N\) عن \(N_{\text{opt}}\)، تزداد الأخطاء، وعندما تبلغ عتبة حرجة \(\varepsilon_{\text{crit}}\)، يخضع النظام لتحول طوري (مثل الانشطار إلى مجموعات فرعية). وبالتالي
	\begin{equation}
		\varepsilon(N_{\text{crit}}) = \varepsilon_{\text{crit}}.
	\end{equation}
	إذا كانت أحجام المجموعات الفرعية بعد التحول قريبة من الحجم الأمثل، فبمبدأ حفظ العدد \(N_{\text{crit}} \approx k N_{\text{opt}}\) مع \(k > 1\). تعتمد القيمة الدقيقة لـ \(k\) على طبيعة التحول (انقسام متماثل أو غير متماثل). باستخدام (\ref{eq:scaling})، نحصل على
	\begin{equation}
		k^{\beta} = \frac{\varepsilon_{\text{crit}}}{\varepsilon_{\text{opt}}}, 
		\qquad\text{وبالتالي}\qquad 
		k = \left( \frac{\varepsilon_{\text{crit}}}{\varepsilon_{\text{opt}}} \right)^{1/\beta}.
		\label{eq:k}
	\end{equation}
	وبالتالي، فإن النسبة \(k\) تتحدد فقط بنسبة الخطأ الحرج إلى الخطأ الأمثل، مرفوعة إلى الأس \(1/\beta\). إذا كانت \(\varepsilon_{\text{crit}}/\varepsilon_{\text{opt}}\) ثابتة تقريبًا عبر الأنظمة المختلفة (وهي فرضية معقولة إذا كانت العتبة محددة بقيد فيزيائي أو بيولوجي أساسي)، فيجب أن تكون \(k\) عالمية أيضًا، ويمكن مقارنة قيمتها بالبيانات التجريبية.
	
	\section{البيانات التجريبية حول الأحجام الحرجة للمجموعات}
	\label{sec:data}
	
	قمنا بجمع بيانات من عدة أنظمة مدروسة جيدًا حيث تم توثيق أحجام المجموعات ونقاط التحول.
	
	\subsection{مستعمرات نحل العسل (Apis mellifera)}
	\label{subsec:bees}
	
	توفر ممارسة تربية النحل أرقامًا موثوقة:
	\begin{itemize}
		\item الحجم الأمثل للمستعمرة (السرب الأول، ``السرب الرئيسي''): \(N_{\text{opt}} = 50\,000\)–\(60\,000\) عاملة (متوسط الكتلة 7–8 كجم، كتلة الفرد \(\approx 0.12\) جم). \cite{Seeley2010, Winston1987}
		\item حجم المستعمرة قبل التطريد: \(N_{\text{crit}} = 80\,000\)–\(100\,000\) (وأحيانًا أكثر). \cite{bees}
		\item وبالتالي \(k_{\text{bees}} = N_{\text{crit}}/N_{\text{opt}} \approx 1.5\)–\(1.67\).
	\end{itemize}
	يكون السرب غير متماثل: يأخذ السرب الرئيسي الملكة القديمة ويترك مستعمرة أصغر. تبلغ \(k\) الملاحظة حوالي 1.6.
	
	\subsection{الوحدات العسكرية}
	\label{subsec:military}
	
	يظهر التنظيم العسكري الحديث هيكلًا هرميًا واضحًا بأحجام وحدات محددة جيدًا \cite{FM3-90.5}:
	\begin{itemize}
		\item الفريق (Squad): 8–12 جنديًا
		\item الفصيلة (Platoon): 30–50 جنديًا
		\item السرية (Company): 100–200 جندي
		\item الكتيبة (Battalion): 300–1000 جندي
		\item اللواء (Brigade): 3000–5000 جندي
		\item الفرقة (Division): 9000–18000 جندي
	\end{itemize}
	تتراوح النسب بين المستويات المتتالية من 3 إلى 5، متجمعة حول 3–4. بالنسبة للسرية، عندما يتجاوز حجمها ~200، قد تنقسم إلى فصيلتين. وهذا يعطي \(k \approx 2\)، وهو أقل من النسبة بين المستويات. هناك حاجة إلى بيانات أفضل عن أحداث الانشطار الفعلية، لكن وجود مستويات هرمية واضحة بحد ذاته يشير إلى قفزات كمية في التعقيد التنظيمي.
	
	\subsection{أسراب الأسماك – سمك الزيبرا (Danio rerio)}
	\label{subsec:fish}
	
	تُظهر تجارب مضبوطة على السلوك الجماعي لسمك الزيبرا \cite{Tunstrom2013} تحولًا طوريًا من التمدرس المنتظم إلى الدوران غير المنتظم مع زيادة الكثافة. تقابل الكثافة الحرجة حجم مجموعة يتراوح بين 50–100 سمكة في حوض قياسي. لم يُعطَ \(N_{\text{opt}}\) دقيق، لكن التحول يحدث عندما يتجاوز المجموعة حجمًا معينًا. وهذا يتسق مع فكرة أن \(k\) حوالي 2.
	
	\subsection{مجموعات الدلافين (Tursiops spp.)}
	\label{subsec:dolphins}
	
	تذكر الدراسات الميدانية للدلافين قارورية الأنف \cite{Connor2000} أحجامًا نموذجية للمجموعات تتراوح بين 30–70 فردًا. في وجود مفترسات (أسماك القرش)، قد تتجمع المجموعات في قطعان تصل إلى عدة مئات (حتى 600). هذا التجمع هو تحول دفاعي، وليس انشطارًا، لكنه يوضح تغيرًا حرجًا في التنسيق. تبلغ نسبة القطيع الكبير إلى المجموعة النموذجية حوالي 5–10، وهو ما قد يقابل \(k\) في الاتجاه المعاكس (الاندماج).
	
	\subsection{تجمعات الماشية – الحفظ الوراثي}
	\label{subsec:livestock}
	
	بيانات من منظمة الأغذية والزراعة (FAO) حول الأحجام الحرجة للسكان للحفظ الوراثي \cite{FAO2013}:
	\begin{itemize}
		\item الماشية: حالة حرجة عندما يكون عدد السكان <2000، معرضة للخطر عند 5000–15000، طبيعي >25000.
		\item الأغنام: يُوصى بـ 100–250 نعجة لحفظ التجمع الجيني.
		\item الخنازير: 25 ذكرًا + 100 أنثى.
		\item الدجاج: 50 ديكًا + 250 دجاجة.
	\end{itemize}
	تعكس هذه الأرقام الحد الأدنى لحجم السكان القابل للحياة، والذي يمكن اعتباره عتبة حرجة للتنسيق الوراثي (تجنب انخفاض التزاوج الداخلي). بتفسير \(N_{\text{opt}}\) كحجم الصحة الوراثية المثلى (مثل 25000 للماشية) و\(N_{\text{crit}}\) كالحد الأدنى القابل للحياة (مثل 2000)، نحصل على \(k \approx 0.08\)، وهو أقل من 1 – وهذا نوع مختلف من التحول (الانقراض). ومع ذلك، قد تتبع النسبة أيضًا قياسًا مع \(\beta\) بطريقة ما، لكنها تتطلب نموذجًا منفصلاً.
	
	\subsection{مجموعات الرئيسيات}
	\label{subsec:primates}
	
	تُظهر بيانات عن قردة البابون \cite{Alberts1995} أحجامًا نموذجية تتراوح بين 30–80 فردًا. عندما تصبح المجموعة كبيرة جدًا (أكثر من ~100)، قد تنقسم إلى مجموعتين ابنتين. وهذا يعطي \(k \approx 2\)–\(3\).
	
	\section{المقارنة مع تنبؤ YUCT}
	\label{sec:comparison}
	
	وفقًا للمعادلة (\ref{eq:k})، تعتمد النسبة \(k = N_{\text{crit}}/N_{\text{opt}}\) على نسبة الخطأ \(\varepsilon_{\text{crit}}/\varepsilon_{\text{opt}}\). إذا كانت هذه النسبة عالمية، فيجب أن تكون \(k\) ثابتة عبر الأنواع. تشير البيانات المجمعة في الجدول \ref{tab:summary} إلى أن \(k\) هي بالفعل في نطاق ضيق نسبيًا، حوالي 1.5–2.5.
	
	\begin{table}[htbp]
		\centering
		\caption{ملخص لنسب الحجم الحرج لمختلف الأنظمة}
		\label{tab:summary}
		\begin{tabular}{lccc}
			\toprule
			\textbf{النظام} & \(N_{\text{opt}}\) & \(N_{\text{crit}}\) & \(k = N_{\text{crit}}/N_{\text{opt}}\) \\
			\midrule
			مستعمرة نحل العسل & 50–60 ألف & 80–100 ألف & 1.5–1.7 \\
			وحدة عسكرية (فصيلة ← سرية) & 30–50 & 100–200 & 2–4 \\
			سرب سمك الزيبرا & 50? & 100? & 2? \\
			مجموعة بابون & 30–80 & ~100 & 1.3–3 \\
			\bottomrule
		\end{tabular}
	\end{table}
	
	إذا أخذنا \(k \approx 1.6\) كقيمة تمثيلية (النحل، الرئيسيات)، فمن المعادلة (\ref{eq:k}) نحصل على
	\[
	\frac{\varepsilon_{\text{crit}}}{\varepsilon_{\text{opt}}} = k^{\beta} = 1.6^{0.67} \approx 1.37.
	\]
	هذا يعني أن الخطأ الحرج أكبر بنسبة 37% فقط من الخطأ الأمثل – وهي عتبة معقولة لبدء تحول طوري. بالنسبة لـ \(k \approx 2\) (الأسماك، بعض مجموعات الرئيسيات)، نحصل على \(2^{0.67} \approx 1.59\). يشير اتساق هذه القيم إلى أن \(\varepsilon_{\text{crit}}/\varepsilon_{\text{opt}}\) ثابت تقريبًا عبر الأنظمة، مما يدعم عالمية \(\beta\) وتفسير العتبة التكنولوجية/التنظيمية كتحول طوري في كفاءة التنسيق.
	
	\section{تنبؤات قابلة للاختبار}
	\label{sec:predictions}
	
	بناءً على الإطار أعلاه، نصيغ عدة تنبؤات محددة يمكن اختبارها تجريبيًا أو رصديًا.
	
	\begin{enumerate}
		\item \textbf{نحل العسل:} يجب أن تكون نسبة حجم السرب الرئيسي إلى حجم المستعمرة قبل التطريد ثابتة عبر السلالات والظروف البيئية المختلفة، ويجب أن تحقق \(k = (\varepsilon_{\text{crit}}/\varepsilon_{\text{opt}})^{1/\beta}\) مع \(\beta = 0.67\). إذا أمكن تقدير \(\varepsilon_{\text{crit}}/\varepsilon_{\text{opt}}\) بشكل مستقل (مثلًا من كفاءة البحث عن الطعام أو معدل الإصابة بالأمراض)، فيجب أن تتطابق \(k\) المتوقعة مع الملاحظات.
		
		\item \textbf{أسراب الأسماك:} في التجارب المضبوطة مع سمك الزيبرا أو الأنواع الأخرى، يجب أن يتناسب حجم المجموعة الذي يحدث عنده التحول من التمدرس إلى الدوران (أو الانشطار) مع الحجم الأمثل وفقًا للقانون نفسه. يمكن أن توفر قياسات أخطاء التنسيق (مثل تباين اتجاه الحركة) تقديرات مستقلة لـ \(\varepsilon_{\text{opt}}\) و \(\varepsilon_{\text{crit}}\).
		
		\item \textbf{الرئيسيات:} يجب أن تسجل الدراسات الميدانية طويلة الأمد لمجموعات البابون والمكاك والشمبانزي أحداث الانشطار وأحجام المجموعات الأم والابنة. يجب أن تكون النسبة \(k\) مستقرة ومتسقة مع تنبؤ YUCT.
		
		\item \textbf{المنظمات العسكرية:} يمكن فحص البيانات التاريخية عن انقسام الوحدات العسكرية (مثل انقسام كتيبة إلى اثنتين). يجب أن تتبع الأحجام التي تحدث عندها هذه الانقسامات نفس القياس، ربما مع \(k\) مختلف بسبب الطبيعة الهرمية.
		
		\item \textbf{تجمعات الماشية:} بالنسبة للحفظ الوراثي، قد يتعلق حجم السكان الفعال \(N_e\) اللازم للحفاظ على التنوع الجيني بالحجم الحرج للانقراض. قد تتناسب نسبة \(N_e\) الأدنى القابل للحياة إلى \(N_e\) الأمثل للبقاء طويل الأمد أيضًا مع \(\beta\)، لكن هذا أكثر تخمينية.
		
		\item \textbf{الشبكات الاجتماعية:} غالبًا ما تنقسم المجموعات الاجتماعية عبر الإنترنت (مثل محادثات واتساب، مجتمعات ريديت) عندما تتجاوز حجمًا معينًا. يمكن تحليل الحجم الحرج للانقسام والحجم النموذجي للمجموعات الابنة باستخدام البصمات الرقمية، مما يوفر مجموعة بيانات ضخمة لاختبار عالمية \(k\) و \(\beta\).
	\end{enumerate}
	
	\section{الارتباط ببديهية العتبة التكنولوجية (الملحق O)}
	\label{sec:connection}
	
	تنص بديهية عتبة القاموس التكنولوجي (الملحق O) على أن النظام يجب أن يصل إلى مستوى تكنولوجي أدنى \(\Tmin\) ليتمكن من استخدام أو توليد قاموس. في سياق تحولات حجم المجموعة، يمكن إعادة تفسير \(\Tmin\) على أنه كفاءة التنسيق الحرجة \(K_{\mathrm{eff},\text{crit}}\) المطلوبة للحفاظ على التماسك. تصبح النسبة التجريبية \(k = N_{\text{crit}}/N_{\text{opt}}\) مقياسًا لمدى نمو النظام قبل أن يحتاج إلى الانشطار أو إعادة التنظيم. حقيقة أن \(k\) من رتبة 1.5–2 عبر أنظمة متنوعة تشير إلى أن النسبة الأساسية \(K_{\mathrm{eff},\text{crit}}/K_{\mathrm{eff},\text{opt}}\) عالمية، بما يتفق مع قانون القياس الكسيري. وبالتالي، تقدم البيانات المعروضة هنا دعمًا غير مباشر لوجود عتبة تكنولوجية وارتباطها بالأس العالمي \(\beta\).
	
	\section{الخلاصة}
	\label{sec:conclusion}
	
	تشير البيانات المجمعة في هذا الملحق، وإن كانت أولية، إلى أن الأحجام الحرجة للمجموعات التي تخضع عندها الأنظمة الجماعية لتحولات طورية ليست عشوائية بل تتبع قانون قياس عالمي مرتبط بأس YUCT \(\beta = 0.67\) وبديهية عتبة القاموس التكنولوجي. تتجمع النسبة \(k = N_{\text{crit}}/N_{\text{opt}}\) حول 1.5–2.5، وعتبة الخطأ الضمنية \(\varepsilon_{\text{crit}}/\varepsilon_{\text{opt}} \approx 1.3\)–1.6 متسقة بشكل ملحوظ عبر أنظمة متنوعة. لقد حددنا تنبؤات محددة قابلة للاختبار يمكن التحقق منها من خلال تجارب ودراسات ميدانية مستقبلية. من شأن تأكيد هذه التنبؤات أن يوفر دعمًا مستقلاً قويًا لـ YUCT وللطبيعة الأساسية لمقياس كفاءة التنسيق \(\Keff\) وقوانين قياسه.
	
	\begin{thebibliography}{99}
		\bibitem{Seeley2010} Seeley, T. D. (2010). \emph{Honeybee Democracy}. Princeton University Press.
		\bibitem{Winston1987} Winston, M. L. (1987). \emph{The Biology of the Honey Bee}. Harvard University Press.
		\bibitem{FM3-90.5} FM 3-90.5 (2008). \emph{Combined Arms Battalion}. US Army Field Manual.
		\bibitem{Tunstrom2013} Tunstrøm, K. et al. (2013). Collective states, multistability and transitional behavior in schooling fish. \emph{PLoS Computational Biology}, 9(2), e1002915.
		\bibitem{Connor2000} Connor, R. C. et al. (2000). The bottlenose dolphin: social relationships in a fission‑fusion society. In: \emph{Cetacean Societies}, University of Chicago Press.
		\bibitem{FAO2013} FAO (2013). \emph{In vivo conservation of animal genetic resources}. FAO Animal Production and Health Guidelines No. 14.
		\bibitem{Alberts1995} Alberts, S. C. \& Altmann, J. (1995). Balancing costs and opportunities: dispersal in male baboons. \emph{American Naturalist}, 145(2), 279–306.
		\bibitem{Yakushev2026} Yakushev, A. V. (2026). \emph{Yakushev Unified Coordination Theory (YUCT)}. Zenodo. \url{https://doi.org/10.5281/zenodo.18444599}
		\bibitem{AppendixL} الملحق L: قياس خطأ التنسيق الكسيري – قانون عالمي من الحمض النووي إلى علم الكون.
		\bibitem{AppendixO} الملحق O: بديهية عتبة القاموس التكنولوجي.
		\bibitem{bees} Seeley, T. D. (2010) – نفس المرجع أعلاه، أو يمكن استبداله بمرجع محدد.
	\end{thebibliography}
	
\end{document}